\chapter{隐马尔可夫模型}


隐马尔可夫模型(Hidden Markov Model, HMM)是自然语言处理中最经典的概率图模型之一.
它常用于词性标注、分词、语音识别等序列建模任务.

HMM的核心思想是: 序列中的状态是不可直接观察的, 但每个隐藏状态会以一定概率生成一个可观测符号.
因此, 我们能够通过观测序列去反推最可能的隐藏状态序列.

\section{贝叶斯定理}


贝叶斯定理给出了先验概率(prior probability)与后验概率(posterior probability)之间的关系:
\[
P(A|B) = (P(B|A) P(A)) / P(B)
\]

其中:

\begin{itemize}
\item $P(A)$是事件$A$的先验概率.
\item $P(B|A)$是在$A$发生条件下$B$发生的概率.
\item $P(A|B)$是在观察到$B$之后对$A$的后验概率.
\end{itemize}

若$A_1, A_2, \dots, A_n$两两互斥且构成完备事件组, 则有
\[
P(A_i|B) = (P(B|A_i) P(A_i)) / (\sum_j P(B|A_j) P(A_j))
\]

贝叶斯定理的重要意义在于: 当直接求后验概率困难时, 可以通过似然与先验来间接计算.
HMM中对状态的推断、参数学习与解码都体现了这种“由观测反推隐藏状态”的思想.

\section{马尔可夫模型的定义}


马尔可夫链(Markov chain)描述的是一个状态序列$Q = (q_1, q_2, \dots, q_T)$的随机演化过程.
若该序列满足一阶马尔可夫性质:
\[
P(q_t | q_1, q_2, \dots, q_{t-1}) = P(q_t | q_{t-1})
\]
则称当前时刻的状态只依赖于前一时刻的状态, 而与更早历史无关.

若状态集合为$\{s_1, s_2, \dots, s_N\}$, 则状态转移概率可写为:
\[
a_{i j} = P(q_{t+1} = s_j | q_t = s_i)
\]
从而得到状态转移矩阵
\[
A = [a_{i j}]_{N \times N}
\]

仅有状态转移而没有观测输出时, 这还是普通马尔可夫链.
HMM在此基础上进一步引入了“观测由状态产生”这一层.

\subsection{三个盒子的直观例子}


假设有$3$个盒子, 每个盒子里分别装有红球和白球:

\begin{table}[H]
\centering
\begingroup
\setlength{\tabcolsep}{2pt}
\small
\begin{tabular}{|p{0.20\textwidth}|p{0.20\textwidth}|p{0.20\textwidth}|p{0.20\textwidth}|}
\hline
 & 盒子1 & 盒子2 & 盒子3 \\ \hline
红球 & 5 & 4 & 7 \\ \hline
白球 & 5 & 6 & 3 \\ \hline
\end{tabular}
\endgroup
\end{table}
\begin{table}[H]
\centering
\begingroup
\setlength{\tabcolsep}{2pt}
\small
\begin{tabular}{|p{0.20\textwidth}|p{0.20\textwidth}|p{0.20\textwidth}|p{0.20\textwidth}|}
\hline
 & 盒子1 & 盒子2 & 盒子3 \\ \hline
抽到红球的概率 & 0.5 & 0.4 & 0.7 \\ \hline
抽到白球的概率 & 0.5 & 0.6 & 0.3 \\ \hline
\end{tabular}
\endgroup
\end{table}
\begin{table}[H]
\centering
\begingroup
\setlength{\tabcolsep}{2pt}
\small
\begin{tabular}{|p{0.20\textwidth}|p{0.20\textwidth}|p{0.20\textwidth}|p{0.20\textwidth}|}
\hline
 & 盒子1 & 盒子2 & 盒子3 \\ \hline
初始选择概率 & 0.2 & 0.4 & 0.4 \\ \hline
\end{tabular}
\endgroup
\end{table}

再假设每次抽完球之后, 下一时刻会按照固定概率切换到另一个盒子.
状态转移矩阵给定为:
\[
A = \begin{bmatrix}0.5 & 0.2 & 0.3 \\ 0.3 & 0.5 & 0.2 \\ 0.2 & 0.3 & 0.5\end{bmatrix}
\]

在这个例子中:

\begin{itemize}
\item 隐状态(hidden state)是“当前选择的是哪个盒子”.
\item 观测(observation)是“当前抽到的球的颜色”.
\item 盒子是看不见的, 球的颜色是看得见的.
\end{itemize}

因此, 它正好构成一个 HMM:

\begin{itemize}
\item 状态集合$S = \{s_1, s_2, s_3\}$.
\item 观测集合$V = \{v_1, v_2\}$, 其中$v_1$表示红球, $v_2$表示白球.
\item 初始概率向量$\pi = [0.2, 0.4, 0.4]$.
\item 状态转移矩阵$A$如上.
\item 发射概率矩阵
\end{itemize}
\[
B = \begin{bmatrix}0.5 & 0.5 \\ 0.4 & 0.6 \\ 0.7 & 0.3\end{bmatrix}
\]

这里$B$的第$i$行表示状态$s_i$下生成不同观测的概率.
例如, 当处于盒子$3$时, 抽到红球的概率为$0.7$, 抽到白球的概率为$0.3$.
$B$的第$i$列则是不同状态观测到同样结果的概率.
例如, 向量$\begin{bmatrix}0.5 \\ 0.6 \\ 0.3\end{bmatrix}$表示各个盒子抽到白球的概率.

\subsection{隐马尔可夫模型的形式化定义}


一个 HMM 通常由三部分参数决定:
\[
\lambda = (A, B, \pi)
\]
其中:

\begin{itemize}
\item $A = [a_{i j}]$是状态转移概率矩阵.
\item $B = [b_j (k)]$是观测概率矩阵, 其中$b_j (k) = P(o_t = v_k | q_t = s_j)$.
\item $\pi = [\pi_i]$是初始状态概率向量, 其中$\pi_i = P(q_1 = s_i)$.
\end{itemize}

设状态序列为
\[
I = (i_1, i_2, \dots, i_T)
\]
观测序列为
\[
O = (o_1, o_2, \dots, o_T)
\]
则 HMM 的两个基本假设是:

\begin{itemize}
\item 一阶马尔可夫假设: 当前状态只依赖前一状态.
\item 观测独立性假设: 当前观测只依赖当前状态.
\end{itemize}

在这两个假设下, 状态序列与观测序列的联合概率为:
\[
P(I, O | \lambda)
    = \pi_{i_1} b_{i_1}(o_1) \prod_{t=2}^T a_{i_{t-1} i_t} b_{i_t}(o_t)
\]

观测序列的概率则可表示为对所有可能状态序列求和:
\[
P(O | \lambda) = \sum_I P(I, O | \lambda)
\]

这个公式是 HMM 三个基本问题的出发点.

\section{马尔可夫模型的三个基本问题}


HMM 的应用通常围绕以下三个基本问题展开:

\begin{itemize}
\item 识别问题(evaluation): 给定模型和观测序列, 计算观测序列出现的概率.
\item 学习问题(learning): 给定观测序列, 估计模型参数.
\item 解码问题(decoding): 给定模型和观测序列, 求最可能的隐藏状态序列.
\end{itemize}

\subsection{识别问题}


已知
\[
\lambda = (A, B, \pi), \quad O = (o_1, o_2, \dots, o_T)
\]
求
\[
P(O | \lambda)
\]

如果直接枚举所有状态序列, 复杂度是指数级的.
因此通常采用动态规划, 即前向算法与后向算法.

\subsubsection{前向算法}


定义前向概率:
\[
\alpha_t (i) = P(o_1, o_2, \dots, o_t, q_t = s_i | \lambda)
\]
它表示: 在时刻$t$处于状态$s_i$, 且前$t$个观测已经出现的联合概率.

前向算法分三步:

\begin{enumerate}
\item[1.] 初始化
\end{enumerate}
\[
\alpha_1 =
\begin{bmatrix}\alpha_1(1) \\ \alpha_1(2) \\ \vdots \\ \alpha_1(N)\end{bmatrix}
= \operatorname{diag}(\pi) b(o_1)
\]

\begin{enumerate}
\item[2.] 递推
\end{enumerate}
\[
\alpha_{t+1} = \operatorname{diag}[b(o_{t+1})] A^T \alpha_t
\]

\begin{enumerate}
\item[3.] 终止
\end{enumerate}
\[
P(O | \lambda) = \sum_i \alpha_T (i)
\]

前向算法将原本指数级的求和压缩为$O(N^2 T)$的计算.

\subsubsection{后向算法}


定义后向概率:
\[
\beta_t (i) = P(o_{t+1}, o_{t+2}, \dots, o_T | q_t = s_i, \lambda)
\]
它表示: 已知时刻$t$处于状态$s_i$, 从$t+1$时刻到终点的剩余观测序列出现的概率.

后向算法也分三步:

\begin{enumerate}
\item[1.] 初始化
\end{enumerate}
\[
\beta_T = \begin{bmatrix}1 \\ 1 \\ \vdots \\ 1\end{bmatrix}
\]

\begin{enumerate}
\item[2.] 递推
\end{enumerate}
\[
\beta_t = A \operatorname{diag}[b(o_{t+1})]\beta_{t+1}
\]

\begin{enumerate}
\item[3.] 终止
\end{enumerate}
\[
P(O | \lambda) = \pi^T \operatorname{diag}[b(o_1)]\beta_1
\]

前向算法与后向算法给出的$P(O | \lambda)$应当一致.
二者也为后续的参数学习算法提供了中间量.

我们以开头的例子来演示前向后向算法如何计算. 假设观测到
\[
O = (\text{红}, \text{白}, \text{红})
\]
记
\[
b(\text{红}) = \begin{bmatrix}0.5 \\ 0.4 \\ 0.7\end{bmatrix}, \quad b(\text{白}) = \begin{bmatrix}0.5 \\ 0.6 \\ 0.3\end{bmatrix}
\]
对应的对角矩阵为
\[
\operatorname{diag}[b(\text{红})] = \begin{bmatrix}0.5 & 0 & 0 \\ 0 & 0.4 & 0 \\ 0 & 0 & 0.7\end{bmatrix}
\]
\[
\operatorname{diag}[b(\text{白})] = \begin{bmatrix}0.5 & 0 & 0 \\ 0 & 0.6 & 0 \\ 0 & 0 & 0.3\end{bmatrix}
\]

先计算前向概率

\begin{enumerate}
\item[1.] 初始化
\end{enumerate}
\[
\alpha_1 = \operatorname{diag}(\pi) b(\text{红})
\]
\[
=
  \begin{bmatrix}0.2 & 0 & 0 \\ 0 & 0.4 & 0 \\ 0 & 0 & 0.4\end{bmatrix}
  \begin{bmatrix}0.5 \\ 0.4 \\ 0.7\end{bmatrix}
\]
\[
= \begin{bmatrix}0.10 \\ 0.16 \\ 0.28\end{bmatrix}
\]

\begin{enumerate}
\item[2.] 递推到$t=2$
\end{enumerate}
\[
\alpha_2 = \operatorname{diag}[b(\text{白})] A^T \alpha_1
\]
先算
\[
A^T \alpha_1
    =
    \begin{bmatrix}0.5 & 0.3 & 0.2 \\ 0.2 & 0.5 & 0.3 \\ 0.3 & 0.2 & 0.5\end{bmatrix}
    \begin{bmatrix}0.10 \\ 0.16 \\ 0.28\end{bmatrix}
\]
\[
= \begin{bmatrix}0.154 \\ 0.184 \\ 0.202\end{bmatrix}
\]
再左乘发射概率对角矩阵:
\[
\alpha_2
    =
    \begin{bmatrix}0.5 & 0 & 0 \\ 0 & 0.6 & 0 \\ 0 & 0 & 0.3\end{bmatrix}
    \begin{bmatrix}0.154 \\ 0.184 \\ 0.202\end{bmatrix}
\]
\[
= \begin{bmatrix}0.077 \\ 0.1104 \\ 0.0606\end{bmatrix}
\]

\begin{enumerate}
\item[3.] 递推到$t=3$
\end{enumerate}
\[
\alpha_3 = \operatorname{diag}[b(\text{红})] A^T \alpha_2
\]
先算
\[
A^T \alpha_2
    =
    \begin{bmatrix}0.5 & 0.3 & 0.2 \\ 0.2 & 0.5 & 0.3 \\ 0.3 & 0.2 & 0.5\end{bmatrix}
    \begin{bmatrix}0.077 \\ 0.1104 \\ 0.0606\end{bmatrix}
\]
\[
= \begin{bmatrix}0.08374 \\ 0.08878 \\ 0.07548\end{bmatrix}
\]
再左乘发射概率对角矩阵:
\[
\alpha_3
    =
    \begin{bmatrix}0.5 & 0 & 0 \\ 0 & 0.4 & 0 \\ 0 & 0 & 0.7\end{bmatrix}
    \begin{bmatrix}0.08374 \\ 0.08878 \\ 0.07548\end{bmatrix}
\]
\[
= \begin{bmatrix}0.04187 \\ 0.035512 \\ 0.052836\end{bmatrix}
\]

\begin{enumerate}
\item[4.] 终止
\end{enumerate}
\[
P(O | \lambda)
    = 0.04187 + 0.035512 + 0.052836
    = 0.130218
\]

再计算后向概率.

\begin{enumerate}
\item[1.] 初始化
\end{enumerate}
\[
\beta_3 = \begin{bmatrix}1 \\ 1 \\ 1\end{bmatrix}
\]

\begin{enumerate}
\item[2.] 递推到$t=2$
\end{enumerate}
\[
\beta_2 = A \operatorname{diag}[b(\text{红})] \beta_3
\]
先算
\[
\operatorname{diag}[b(\text{红})] \beta_3
    =
    \begin{bmatrix}0.5 & 0 & 0 \\ 0 & 0.4 & 0 \\ 0 & 0 & 0.7\end{bmatrix}
    \begin{bmatrix}1 \\ 1 \\ 1\end{bmatrix}
\]
\[
= \begin{bmatrix}0.5 \\ 0.4 \\ 0.7\end{bmatrix}
\]
再左乘状态转移矩阵:
\[
\beta_2
    =
    \begin{bmatrix}0.5 & 0.2 & 0.3 \\ 0.3 & 0.5 & 0.2 \\ 0.2 & 0.3 & 0.5\end{bmatrix}
    \begin{bmatrix}0.5 \\ 0.4 \\ 0.7\end{bmatrix}
\]
\[
= \begin{bmatrix}0.54 \\ 0.49 \\ 0.57\end{bmatrix}
\]

\begin{enumerate}
\item[3.] 递推到$t=1$
\end{enumerate}
\[
\beta_1 = A \operatorname{diag}[b(\text{白})] \beta_2
\]
先算
\[
\operatorname{diag}[b(\text{白})] \beta_2
    =
    \begin{bmatrix}0.5 & 0 & 0 \\ 0 & 0.6 & 0 \\ 0 & 0 & 0.3\end{bmatrix}
    \begin{bmatrix}0.54 \\ 0.49 \\ 0.57\end{bmatrix}
\]
\[
= \begin{bmatrix}0.27 \\ 0.294 \\ 0.171\end{bmatrix}
\]
再左乘状态转移矩阵:
\[
\beta_1
    =
    \begin{bmatrix}0.5 & 0.2 & 0.3 \\ 0.3 & 0.5 & 0.2 \\ 0.2 & 0.3 & 0.5\end{bmatrix}
    \begin{bmatrix}0.27 \\ 0.294 \\ 0.171\end{bmatrix}
\]
\[
= \begin{bmatrix}0.2451 \\ 0.2622 \\ 0.2277\end{bmatrix}
\]

\begin{enumerate}
\item[4.] 终止
\end{enumerate}
\[
P(O | \lambda) = \pi^T \operatorname{diag}[b(\text{红})] \beta_1
\]
先算
\[
\operatorname{diag}[b(\text{红})] \beta_1
    =
    \begin{bmatrix}0.5 & 0 & 0 \\ 0 & 0.4 & 0 \\ 0 & 0 & 0.7\end{bmatrix}
    \begin{bmatrix}0.2451 \\ 0.2622 \\ 0.2277\end{bmatrix}
\]
\[
= \begin{bmatrix}0.12255 \\ 0.10488 \\ 0.15939\end{bmatrix}
\]
再与$\pi^T$相乘:
\[
P(O | \lambda)
    = [0.2, 0.4, 0.4] \begin{bmatrix}0.12255 \\ 0.10488 \\ 0.15939\end{bmatrix}
\]
\[
= 0.02451 + 0.041952 + 0.063756
\]
\[
= 0.130218
\]

因此, 按上述矩阵递推可以把前向变量与后向变量逐步算出.
并且前向算法与后向算法都得到相同的结果$0.130218$.

\subsection{解码问题}


已知
\[
\lambda = (A, B, \pi), \quad O = (o_1, o_2, \dots, o_T)
\]
求使得
\[
P(I | O, \lambda)
\]
最大的状态序列
\[
I = (i_1, i_2, \dots, i_T)
\]

这就是寻找最可能隐藏路径的问题.
最常用的方法是维特比算法(Viterbi algorithm).

\subsubsection{维特比算法}


定义
\[
\delta_t (i) = \max_{i_1, \dots, i_{t-1}} P(i_1, \dots, i_t = s_i, o_1, \dots, o_t | \lambda)
\]
它表示: 在时刻$t$到达状态$s_i$的所有路径中, 概率最大的那条路径的概率.

同时定义路径回溯指针:
\[
\psi_t(j) = \arg \max_i [\delta_{t-1}(i) a_{i j}]
\]

维特比算法的步骤如下:

\begin{enumerate}
\item[1.] 初始化
\end{enumerate}
\[
\delta_1 = \operatorname{diag}(\pi) b(o_1), \quad \psi_1 = \begin{bmatrix}0 \\ 0 \\ \vdots \\ 0\end{bmatrix}
\]

\begin{enumerate}
\item[2.] 递推
\end{enumerate}
\[
\delta_{t+1} = \operatorname{diag}[b(o_{t+1})]\operatorname{rowmax} [A^T \operatorname{diag}(\delta_t)]
\]
\[
\psi_{t+1} = \arg \operatorname{rowmax} [A^T \operatorname{diag}(\delta_t)]
\]

\begin{enumerate}
\item[3.] 终止
\end{enumerate}
\[
P^* = \max_i \delta_T (i)
\]
\[
i_T^* = \arg \max_i \delta_T (i)
\]

\begin{enumerate}
\item[4.] 回溯
\end{enumerate}
\[
i_t^* = \psi_{t+1}(i_{t+1}^*), \quad t = T-1, T-2, \dots, 1
\]

与前向算法不同, 维特比算法在每一步保留的不是“所有路径概率之和”, 而是“到达该状态的最大路径概率”.
这正是动态规划思想在最优路径搜索中的应用.

我们仍然用前文那个模型举例说明计算过程.
假设我们观测到序列: $(\text{红}, \text{白}, \text{红}, \text{白})$.

\providecommand{\viterbimax}[1]{\textcolor{red!70!black}{\mathbf{#1}}}

状态顺序仍记为$(s_1, s_2, s_3)$.
这里维特比算法最大化的是$P(I, O | \lambda)$, 由于$P(O | \lambda)$对所有候选路径相同, 它也等价于最大化$P(I | O, \lambda)$.
下面矩阵中红色加粗的数字表示该行取到的最大值.

\begin{enumerate}
\item[1.] 初始化
\end{enumerate}
\[
\delta_1 = \operatorname{diag}(\pi) b(\text{红})
    =
    \begin{bmatrix}0.2 & 0 & 0 \\ 0 & 0.4 & 0 \\ 0 & 0 & 0.4\end{bmatrix}
    \begin{bmatrix}0.5 \\ 0.4 \\ 0.7\end{bmatrix}
    = \begin{bmatrix}0.10 \\ 0.16 \\ 0.28\end{bmatrix}
\]
\[
\psi_1 = \begin{bmatrix}0 \\ 0 \\ 0\end{bmatrix}
\]

\begin{enumerate}
\item[2.] 递推到$t=2$, 观测为白球
\end{enumerate}
\[
A^T \operatorname{diag}(\delta_1)
    =
    \begin{bmatrix}
    0.050 & 0.048 & \viterbimax{0.056} \\
    0.020 & 0.080 & \viterbimax{0.084} \\
    0.030 & 0.032 & \viterbimax{0.140}
    \end{bmatrix}
\]
按行取最大值得到:
\[
\operatorname{rowmax}[A^T \operatorname{diag}(\delta_1)]
    = \begin{bmatrix}\viterbimax{0.056} \\ \viterbimax{0.084} \\ \viterbimax{0.140}\end{bmatrix}
\]
\[
\psi_2 = \begin{bmatrix}3 \\ 3 \\ 3\end{bmatrix}
\]
再乘以当前观测的发射概率:
\[
\delta_2
    = \operatorname{diag}[b(\text{白})] \begin{bmatrix}0.056 \\ 0.084 \\ 0.140\end{bmatrix}
    = \begin{bmatrix}0.0280 \\ 0.0504 \\ 0.0420\end{bmatrix}
\]

\begin{enumerate}
\item[3.] 递推到$t=3$, 观测为红球
\end{enumerate}
\[
A^T \operatorname{diag}(\delta_2)
    =
    \begin{bmatrix}
    0.01400 & \viterbimax{0.01512} & 0.00840 \\
    0.00560 & \viterbimax{0.02520} & 0.01260 \\
    0.00840 & 0.01008 & \viterbimax{0.02100}
    \end{bmatrix}
\]
\[
\operatorname{rowmax}[A^T \operatorname{diag}(\delta_2)]
    = \begin{bmatrix}\viterbimax{0.01512} \\ \viterbimax{0.02520} \\ \viterbimax{0.02100}\end{bmatrix}
\]
\[
\psi_3 = \begin{bmatrix}2 \\ 2 \\ 3\end{bmatrix}
\]
\[
\delta_3
    = \operatorname{diag}[b(\text{红})] \begin{bmatrix}0.01512 \\ 0.02520 \\ 0.02100\end{bmatrix}
    = \begin{bmatrix}0.007560 \\ 0.010080 \\ 0.014700\end{bmatrix}
\]

\begin{enumerate}
\item[4.] 递推到$t=4$, 观测为白球
\end{enumerate}
\[
A^T \operatorname{diag}(\delta_3)
    =
    \begin{bmatrix}
    \viterbimax{0.0037800} & 0.0030240 & 0.0029400 \\
    0.0015120 & \viterbimax{0.0050400} & 0.0044100 \\
    0.0022680 & 0.0020160 & \viterbimax{0.0073500}
    \end{bmatrix}
\]
\[
\operatorname{rowmax}[A^T \operatorname{diag}(\delta_3)]
    = \begin{bmatrix}\viterbimax{0.0037800} \\ \viterbimax{0.0050400} \\ \viterbimax{0.0073500}\end{bmatrix}
\]
\[
\psi_4 = \begin{bmatrix}1 \\ 2 \\ 3\end{bmatrix}
\]
\[
\delta_4
    = \operatorname{diag}[b(\text{白})] \begin{bmatrix}0.0037800 \\ 0.0050400 \\ 0.0073500\end{bmatrix}
    = \begin{bmatrix}0.00189000 \\ 0.00302400 \\ 0.00220500\end{bmatrix}
\]

\begin{enumerate}
\item[5.] 终止与回溯
\end{enumerate}
\[
P^* = \max_i \delta_4(i) = 0.003024, \quad i_4^* = 2
\]
根据回溯指针:
\[
i_3^* = \psi_4(2) = 2, \quad
i_2^* = \psi_3(2) = 2, \quad
i_1^* = \psi_2(2) = 3
\]
因此最可能的隐藏状态序列为:
\[
I^* = (s_3, s_2, s_2, s_2)
\]
对应的路径概率为:
\[
\pi_3 b_3(\text{红}) a_{3 2} b_2(\text{白}) a_{2 2} b_2(\text{红}) a_{2 2} b_2(\text{白})
\]
\[
= 0.4 \times 0.7 \times 0.3 \times 0.6 \times 0.5 \times 0.4 \times 0.5 \times 0.6
  = 0.003024
\]

\subsection{学习问题}


已知
\[
O = (o_1, o_2, \dots, o_T)
\]
求使得
\[
P(O | \lambda)
\]
最大的模型参数
\[
\lambda = (A, B, \pi)
\]

这就是 HMM 的参数估计问题.
当状态序列不可观测时, 通常采用 Baum-Welch 算法, 它本质上是 EM 算法在 HMM 上的具体形式.

\subsubsection{Baum Welch算法}


先定义两个辅助概率.

\begin{enumerate}
\item[1.] 时刻$t$处于状态$s_i$的概率:
\end{enumerate}
\[
\gamma_t (i) = P(q_t = s_i | O, \lambda)
\]
利用前向概率和后向概率可写为:
\[
\gamma_t (i) = (\alpha_t (i) \beta_t (i)) / (\sum_j \alpha_t (j) \beta_t (j))
\]

\begin{enumerate}
\item[2.] 时刻$t$处于$s_i$且时刻$t+1$处于$s_j$的概率:
\end{enumerate}
\[
\xi_t (i, j) = P(q_t = s_i, q_{t+1} = s_j | O, \lambda)
\]
其表达式为:
\[
\xi_t (i, j) =
    (\alpha_t (i) a_{i j} b_j (o_{t+1}) \beta_{t+1}(j))
    /
    (\sum_i \sum_j \alpha_t (i) a_{i j} b_j (o_{t+1}) \beta_{t+1}(j))
\]

据此可以更新参数:

\begin{itemize}
\item 初始概率
\end{itemize}
\[
\pi_i = \gamma_1(i)
\]

\begin{itemize}
\item 状态转移概率
\end{itemize}
\[
a_{i j} = (\sum_{t=1}^{T-1} \xi_t (i, j)) / (\sum_{t=1}^{T-1} \gamma_t (i))
\]

\begin{itemize}
\item 发射概率
\end{itemize}
\[
b_j(k) =
    (\sum_{t:o_t=v_k} \gamma_t (j))
    /
    (\sum_{t=1}^T \gamma_t (j))
\]

不断重复“计算期望”与“更新参数”两个步骤, 就能逐步增大$P(O | \lambda)$, 直到收敛.
